\section{Positive laws of admissible paths}\label{sec:laws}

We now turn failure of a small compatible barrier into a probability law
of paths with controlled mean arclength occupation and boundary contacts.
The construction has two stages.  Modulus duality supplies a law for each
prescribed countable set of dust points, together with an auxiliary
atomless measure on the dust.  A fusion lemma then produces one law with
square-summable hit probabilities on all of $B$.

Two requirements run through the argument.  Admissibility of a metric
must hold for every path, including paths that charge a reference-null set
where the metric is infinite.  The constants in the first stage must also
be uniform in the prescribed countable set and atomless measure.  These
requirements allow the two stages to be combined.

Throughout this section the endpoint data of
Definition~\ref{def:endpoint} are fixed.  The preceding coordinate list is
countable and satisfies \textup{(P)}.  The collection \(\mathscr Q\) of
preceding path laws is countable, and \(m_\xi\in\ell^2(B)\) for every
\(\xi\in\mathscr Q\).  Write
\[
 D_* = \Sigma\cup\bigcup_{\xi\in\mathscr Q}\{s\in B:m_\xi(s)>0\}.
\]
The set $D_*$ is countable; it records the points at which the prior data
will require control of visit multiplicity.  A \emph{base code} is a tagged path
\(\Gamma=(\gamma,\theta)\) satisfying
Definition~\ref{def:path-system}, including
Equations~\eqref{eq:path-law-constraint} and
\eqref{eq:path-level-constraint}, that starts in \(E\), first reaches \(b\)
at time \(1\), and respects the gates when gates are prescribed.  We impose no
auxiliary atomless-measure restriction or tuple exclusions on base codes.
Their space, denoted by
\(\mathscr A\), is standard Borel by Lemma~\ref{lem:path-coding}.

For a code \(\Gamma=(\gamma,\theta)\), write
\(\ell_\Gamma:=\nu_\gamma\) for the arclength occupation measure defined in
Section~\ref{sec:barriers}, and put
\[
 N_\Gamma(s)=\#\{t\in[0,1]:\gamma(t)=s\},\qquad
 N^-_\Gamma(s)=\#\{t\in[0,1):\gamma(t)=s\}.
\]
An infinite cardinality in these expressions means \(+\infty\).
Thus
\(N_\Gamma=N^-_\Gamma+\ind_{\{b\}}\), since the target is first reached at
the final time.  For a law \(\mu\) on base codes, use the notation
\[
 \bar\ell_\mu=\int\ell_\Gamma\,\dd\mu(\Gamma),\qquad
 n_\mu(s)=\E_\mu N_\Gamma(s),\qquad
 m_\mu(s)=\Prob_\mu\{N_\Gamma(s)>0\}.
\]
Here \(n_\mu\) counts visits with parameter multiplicity, whereas \(m_\mu\)
is a hit probability.  In particular, \(m_\mu\le n_\mu\).  We will retain
the stronger visit-count estimate on $D_*$ and obtain the hit-probability
estimate on the entire dust.

\begin{theorem}[Positive-law alternative]\label{thm:positive-law}
Suppose that, for some \(\beta>0\), there is no new compatible barrier
coordinate with
\[
 \|\dd u\|_{2,G\setminus O}+\|a_u\|_{\ell^2(B)}<\beta.
\]
Nonexistence here allows an arbitrary countable enlargement of the flags and,
in the gated problem, any penalty region permitted in
Definition~\ref{def:endpoint}.  Then there is a probability law \(\mu\) on
base codes with the following properties:
\begin{enumerate}[label=\textup{(\roman*)}]
\item \(\bar\ell_\mu=g_\mu\,\dd a\), where
      \(g_\mu\in L^2(G,\dd a)\);
\item \(n_\mu|_{D_*}\in\ell^2(D_*)\) and
      \(m_\mu\in\ell^2(B)\);
\item for every slot and every tag type, its tuple distribution, regarded as
      a subprobability measure, is absolutely continuous with respect to
      Lebesgue measure on the corresponding tuple torus.
\end{enumerate}
The three norms in \textup{(i)--(ii)} are bounded by a constant depending only
on \(\beta\) and the numerical constant in Theorem~\ref{thm:barrier}.
In particular, the bound does not depend on any auxiliary countable subset of
\(B\) or any auxiliary finite atomless measure on \(B\).
\end{theorem}

To prove the theorem, we first establish the measure-family duality used
in Lemma~\ref{lem:countable-law}.  That extraction lemma gives all the
required estimates on prescribed countable data.  We then prove
Lemma~\ref{lem:law-fusion} and check its hypotheses for the resulting
family of path laws.

\subsection{Duality for a Borel family of measures}

Modulus for systems of measures goes back to Fuglede
\cite[Chapter~I]{Fuglede1957}.  Its dual description by probability
measures with controlled barycenter was developed by Ambrosio,
Di Marino, and Savar\'e
\cite[Theorem~5.1 and Corollary~5.2\textup{(b)}]{AmbrosioDiMarinoSavare2015}.
We give the $p=2$ argument for the Borel families needed here.  The input
will be a positive lower bound for modulus; the output will be a law
whose mean measure has a bounded $L^2$ density.  The proof retains
pointwise admissibility for extended densities and produces compact
support after a refinement of the code topology.

If \(Y\) is Polish and \(\sigma\) is a finite Borel measure on \(Y\), let
\(\mathcal M_+(Y)\) denote the space of finite positive Borel measures with
the narrow topology.  We write $\mathcal P(Y)$ for its subspace of
probability measures.  For a family \((\nu_a)_{a\in A}\) of such measures,
define
\[
 \Mod_\sigma(A)=
 \inf\left\{\int_Y h^2\,\dd\sigma:
  h:Y\longrightarrow[0,\infty]\text{ is Borel},\quad
  \int_Y h\,\dd\nu_a\ge1\ \text{for every }a\in A\right\}.
\]
The admissibility condition is pointwise in the index \(a\); no measures
\(\nu_a\) are discarded because of null sets for \(\sigma\).

\begin{lemma}[Measure-family duality]\label{lem:modulus-duality}
Let $Y$ be Polish, let $\sigma$ be a finite Borel measure on $Y$,
let \(A\) be standard Borel, and let
\(a\mapsto\nu_a\in\mathcal M_+(Y)\) be Borel.  If
\(\Mod_\sigma(A)>c>0\), then, for every \(C>c^{-1/2}\), there is a Borel
probability measure \(P\) on \(A\) whose barycenter satisfies
\[
 \int_A\nu_a\,\dd P(a)=q\,\sigma,
 \qquad q\ge0,
 \qquad \|q\|_{L^2(\sigma)}\le C.
\]
Moreover, after giving \(A\) a Polish topology with its original Borel sets
and with \(a\mapsto\nu_a\) continuous, the law may be supported on a compact
subset of \(A\).
\end{lemma}

\begin{proof}
The Polish refinement theorem gives the topology in the statement; see
\cite[Section~13.A, in particular Exercise~13.5]{Kechris1995}.  We work in this topology.  First we prove continuity
under increasing unions and approximation by open neighborhoods, keeping
the prescribed values of every density on reference-null sets.  These two
properties will give a compact subfamily of positive modulus.  Separation
on that compact subfamily will then produce the required probability law.

\smallskip\noindent
\emph{Increasing families.}
If \(A_j\uparrow A_\infty\), then
\begin{equation}\label{eq:modulus-increasing}
 \Mod_\sigma(A_\infty)=\lim_j\Mod_\sigma(A_j).
\end{equation}
Only the case where the limit \(M\) is finite needs proof.  Choose
nonnegative Borel admissible metrics \(h_j\) for \(A_j\) whose squared
norms tend to \(M\).  Weak compactness in \(L^2(\sigma)\), followed by
convex combinations of tails, gives a nonnegative limit \(h\) and finite
convex combinations \(v_k\) of metrics with indices at least \(k\), such
that
\[
 \|h\|_2\le\sqrt M,
 \qquad \sum_k\|v_k-h\|_2<\infty.
\]
Choose a Borel representative of \(h\) finite everywhere.  Keep the actual
extended Borel representatives of the \(v_k\), and set
\(D=\sum_k|v_k-h|\).  This is an extended nonnegative Borel function with
\(\|D\|_2<\infty\).  For \(a\in A_\infty\), the metrics \(v_k\) are
admissible at \(a\) for all sufficiently large \(k\).  If
\(\int D\,\dd\nu_a=\infty\), then \(h+\eps D\) is admissible there for
every \(\eps>0\).  Otherwise
\(\int|v_k-h|\,\dd\nu_a\to0\), so either \(\int h\,\dd\nu_a=\infty\)
or convergence of the integrals gives \(\int h\,\dd\nu_a\ge1\).
Consequently \(h+\eps D\) is admissible on the entire union.  Letting
\(\eps\downarrow0\) proves Equation~\eqref{eq:modulus-increasing}.
The added density $D$ accounts for every index at which convergence in
$L^2(\sigma)$ alone would not control integration against $\nu_a$.

\smallskip\noindent
\emph{Open neighborhoods.}
For every subset \(F\subset A\) and every \(q>\Mod_\sigma(F)\), there is
an open set \(V\supset F\) with \(\Mod_\sigma(V)<q\).
To see this, first note that every nonnegative Borel
\(h\in L^2(\sigma)\), including its prescribed extended values on null
sets, has a nonnegative lower semicontinuous pointwise majorant with
arbitrarily little increase of its \(L^2\) norm.  Indeed, for \(e>0\) set
\[
 h_e=e\sum_{j=0}^{\infty}\ind_{\{h>je\}}.
\]
At finite values, \(h\le h_e\le h+e\).  By outer regularity, choose open
sets \(U_j\supset\{h>je\}\) for which
\(e\sum_j\sigma(U_j\setminus\{h>je\})^{1/2}\) is arbitrarily small.
Then
\(H=e\sum_j\ind_{U_j}\) is lower semicontinuous, majorizes \(h\)
pointwise, and satisfies
\[
 \|H\|_2\le\|h\|_2+e\sigma(Y)^{1/2}
       +e\sum_j\sigma(U_j\setminus\{h>je\})^{1/2}.
\]
Apply this construction to an admissible metric for \(F\) of squared norm
less than \(q\), and scale its majorant slightly so that its integrals over
\(F\) are strictly greater than \(1\), still with squared norm less than
\(q\).  The set
\(V=\{a:\int H\,\dd\nu_a>1\}\) is open by lower semicontinuity of
integration and continuity of the measure map.  It has the required modulus
bound.

\smallskip\noindent
\emph{A compact subfamily.}
We now use the two modulus properties to retain positive modulus while
restricting the codes to a compact set.  Use a complete compatible metric
on \(A\).  At step \(j\), cover \(A\)
by countably many closed balls of radius \(2^{-j}\).  Starting with
\(F_0=A\), Equation~\eqref{eq:modulus-increasing} permits selection of a
finite union of these balls whose intersection \(F_j\) with \(F_{j-1}\)
still has modulus greater than \(c\).  The sets \(F_j\) are nested and
closed.  Every sequence \(a_j\in F_j\) has a Cauchy subsequence, by the
finite ball covers at each fixed depth.  Completeness shows that
\(K=\bigcap_jF_j\) is nonempty and compact.  The same argument shows that
every open neighborhood of \(K\) contains all sufficiently late \(F_j\).
If \(\Mod_\sigma(K)<c\), the open-neighborhood property would therefore
contradict \(\Mod_\sigma(F_j)>c\).  Hence
\(\Mod_\sigma(K)\ge c\).

The compactness here comes from the nested finite covers and completeness.
Individual closed balls in the refined topology need not be compact.

\smallskip\noindent
\emph{Separation and the barycenter.}
For every bounded continuous \(h\ge0\) with \(\|h\|_2\le1\), some
\(a\in K\) satisfies \(\int h\,\dd\nu_a\le C\).  Otherwise \(h/C\)
would be admissible on \(K\), contrary to
\(C^{-2}<c\le\Mod_\sigma(K)\).
Given finitely many such tests \(h_1,\ldots,h_n\), consider their integral
vectors averaged against probability measures on \(K\).  This is a compact
convex subset of \(\R^n\).  If it missed the orthant
\(( -\infty,C]^n\), strict separation would give
\(\alpha_j\ge0\), \(\sum_j\alpha_j=1\), such that
\(\sum_j\alpha_j\int h_j\,\dd\nu_a>C\) for every \(a\in K\).
The test \(\sum_j\alpha_jh_j\) would contradict the preceding observation.
Thus a probability measure meets all the finitely many inequalities.

Compactness of the probability measures on \(K\), and continuity of the
integral tests there, now give one law \(P\) meeting every bounded continuous
nonnegative test.  Its barycenter \(\eta\) is finite: the continuous mass
function \(a\mapsto\nu_a(Y)\) is bounded on \(K\).  Scaling the tests,
including arbitrarily large multiples of tests of zero \(L^2\) norm, gives
\[
 \int h\,\dd\eta\le C\|h\|_{L^2(\sigma)}
 \quad(h\ge0\text{ bounded and continuous}).
\]
Apply this also to \(|h|\).  Density of bounded continuous functions in
\(L^2(\sigma)\) and the Hilbert-space representation theorem give a
nonnegative density \(q\) of norm at most \(C\).  The finite measures
\(\eta\) and \(q\sigma\) agree on bounded continuous functions and hence
are equal.
\end{proof}

\subsection{Countable and atomless tests}

We apply the duality lemma to three measures associated with a tagged
path: its domain arclength occupation, its visits to a prescribed
countable set, and a weighted measure of its tag tuples.  The first two
will give the occupation and visit-count estimates.  The third will make
every slot's tuple law absolutely continuous.  Together with the residual
atomless-measure constraint, this will control the whole path image under
the prescribed atomless measure.

Let \(\mathcal I\) be the countable set of tag types
\(\boldsymbol\iota=(i_1,\ldots,i_d)\), and write \(d(\boldsymbol\iota)=d\).
Fix positive numbers \(\omega_{\boldsymbol\iota}\) with sum at most \(1\).
Use the Polish disjoint union and finite reference measure
\[
 \Pi=\coprod_{\boldsymbol\iota\in\mathcal I}
          \T^{d(\boldsymbol\iota)},
 \qquad
 \tau=\sum_{\boldsymbol\iota\in\mathcal I}
          \omega_{\boldsymbol\iota}\,\dd\mathbf t.
\]
If there are no tag types, \(\Pi\) is empty.  For a code \(\Gamma\), define
its weighted tuple measure by
\begin{equation}\label{eq:weighted-tuple-measure}
 \psi_\Gamma=
 \sum_{k\text{ used}}2^{-k}
       \delta_{(\boldsymbol\iota_k,\mathbf t_k)}.
\end{equation}
The components of \(\Pi\) are indexed by \emph{types}, whereas the factors
\(2^{-k}\) are indexed by optional \emph{slots}.  This distinction allows
tags to be reassigned to later slots without changing their geometric data.

\begin{lemma}[Prescribed-countable-set extraction]\label{lem:countable-law}
Under the hypothesis of Theorem~\ref{thm:positive-law}, there is a constant
\(C<\infty\), depending only on \(\beta\) and the constant in
Theorem~\ref{thm:barrier}, with the following property.  For every countable
\(D\supset D_*\cup\{b\}\) and every finite atomless Borel measure
\(\lambda\) on \(B\), there is a law \(\mu\) on base codes such that
\begin{align}
 \bar\ell_\mu&=g_\mu\,\dd a,
 &\|g_\mu\|_2&\le C,\label{eq:countable-law-occupation}\\
 \|n_\mu|_D\|_{\ell^2(D)}&\le C+1,
 &\int_Bm_\mu\,\dd\lambda&=0,\label{eq:countable-law-dust}\\
 \E_\mu\psi_\Gamma&=q_\mu\,\tau,
 &\|q_\mu\|_{L^2(\tau)}&\le C.\label{eq:countable-law-tuples}
\end{align}
In addition, \(\lambda(\gamma(Z_\Gamma))=0\) for almost every code.
\end{lemma}

\begin{proof}
\smallskip\noindent
\emph{The reference space and path measures.}
Choose \(c_s>0\), \(s\in D\), with \(\sum_{s\in D}c_s^2\le1\), and put
\[
 Y=G\sqcup D_{\mathrm{disc}}\sqcup\Pi,
 \qquad
 \sigma=\dd a\oplus\sum_{s\in D}c_s^2\delta_s\oplus\tau.
\]
Add the constraint \(\lambda(\gamma(Z_\Gamma))=0\) to the base path
system.  For its codes define
\begin{equation}\label{eq:path-coding-measure}
 \nu_\Gamma=\ell_\Gamma\oplus
       \sum_{s\in D}c_sN^-_\Gamma(s)\delta_s\oplus\psi_\Gamma.
\end{equation}
Restrict to codes for which this measure is finite.  This is a standard
Borel family, and its map to \(\mathcal M_+(Y)\) is Borel by
Lemma~\ref{lem:path-coding}.  In particular, the mean measures used below
are ordinary Borel measure kernels.

\smallskip\noindent
\emph{A modulus bound independent of the prescribed tests.}
Let \(C_{\mathrm{bar}}\) denote the numerical constant in
Theorem~\ref{thm:barrier}.  The modulus of this family is at least
\begin{equation}\label{eq:uniform-modulus-lower-bound}
 c_0=\left(\frac{\beta}{4C_{\mathrm{bar}}}\right)^2.
\end{equation}
Suppose, to the contrary, that it has an admissible extended Borel metric
\(h\) with \(\|h\|_{L^2(\sigma)}<\beta/(4C_{\mathrm{bar}})\).  Its
domain and atom parts give
\[
 g_0=h|_G,
 \qquad
 m_0(s)=
 \begin{cases}c_sh(s),&s\in D,\\0,&s\notin D,
 \end{cases}
 \qquad
 \|g_0\|_2^2+\|m_0\|_{\ell^2(B)}^2\le\|h\|_{L^2(\sigma)}^2.
\]
For a sufficiently small \(\delta>0\), set
\(g=g_0+\delta\) and
\(m=m_0+\delta\sum_{s\in D}c_s\ind_{\{s\}}\).  Then \(g\) is bounded
below by a positive constant and
\[
 C_{\mathrm{bar}}\bigl(\|g\|_2+\|m\|_{\ell^2(B)}\bigr)<\beta/2.
\]
The small additions are permitted because \(\area(G)<\infty\) and
\(\|(c_s)\|_{\ell^2(D)}\le1\).

In each type exclude the set on which the corresponding tuple part of
\(h\) is infinite.  It is Lebesgue null, because that component has a
positive reference weight.  We claim that every allowed path for this
additional null-exclusion system has true \((g,m)\)-cost at least \(1\),
omitting the final target toll.  It suffices to consider one of finite true
cost.  The added tolls imply
\(\sum_{s\in D}c_sN^-_\Gamma(s)<\infty\), so the non-tuple parts of
Equation~\eqref{eq:path-coding-measure} are finite.  All the tagged tuples
have finite \(h\)-value.  Reassign the countably many tags to successively
larger unused slots so that the integral of the tuple part of \(h\) is
less than any prescribed \(\eps>0\).  This leaves the path, tags, residual
set, and every path-system constraint unchanged.  The resulting code
belongs to the family on which \(h\) is admissible.  Since its non-tuple
integrals are bounded by the true cost,
\[
 1\le\int h\,\dd\nu_\Gamma
       \le\mathcal L^-_{g,m}(\gamma)+\eps.
\]
Here the superscript minus means that only the final target time is omitted.
Let \(\eps\downarrow0\).  Paths of infinite cost satisfy the same lower
bound automatically.

Theorem~\ref{thm:barrier} applies to this path system, with its atomless
constraint and its additional null tuple exclusions.  Taking the error in
that theorem sufficiently small produces a compatible barrier of budget
less than \(\beta\), a contradiction.  This proves
Equation~\eqref{eq:uniform-modulus-lower-bound}.  Notice that its constant
does not depend on \(D\), \(\lambda\), or the choices of \(c_s\).

\smallskip\noindent
\emph{Occupation, visit counts, and tuple distributions.}
Apply Lemma~\ref{lem:modulus-duality} with \(c=c_0/2\) and any fixed
\(C>\sqrt{2/c_0}\).  The resulting law has mean measure with density of
\(L^2(\sigma)\) norm at most \(C\).  On \(G\) and \(\Pi\), this gives
Equations~\eqref{eq:countable-law-occupation} and
\eqref{eq:countable-law-tuples}.  At \(s\in D\), the mean mass is
\(c_s\E N^-_\Gamma(s)\), while the reference mass is \(c_s^2\).
Consequently the contribution of the atom part to squared density norm is
exactly
\[
 \sum_{s\in D}\bigl(\E N^-_\Gamma(s)\bigr)^2.
\]
Restoring the final visit adds one coordinate vector of norm \(1\), proving
the count estimate in Equation~\eqref{eq:countable-law-dust}.

\smallskip\noindent
\emph{The atomless test on the whole path image.}
It remains to verify the atomless-measure conclusion, which concerns the
whole path image.  Equation~\eqref{eq:countable-law-tuples} implies absolute
continuity of the tuple distribution in every fixed slot and type: that
distribution, multiplied by \(2^{-k}\), is dominated by the mean weighted
tuple measure.  At a fixed free point \(s\), incidence of a tag of type
\((i_1,\ldots,i_d)\) requires at least one of its levels to equal the
unique trace of the corresponding coordinate at \(s\), modulo \(1\).
This is a finite union of Lebesgue-null coordinate sections of the tuple
torus.  Thus a fixed free point is visited during an open tag interval with
probability zero.  Since \(\lambda(\Sigma)=0\), Fubini and the countable
number of slots and types show that the expected \(\lambda\)-measure of
all tagged dust contacts is zero.  The residual contacts have
\(\lambda\)-measure zero on every code in the family.  Combining the two
parts gives \(\int_Bm_\mu\,\dd\lambda=0\).
\end{proof}

\subsection{Fusion from countable tests to the whole dust}

The preceding lemma permits the law to depend on the countable set $D$
and the atomless measure $\lambda$.  We now need a single law whose hit
probabilities satisfy a common bound on every finite subset of $B$.
The next lemma makes this passage using joint upper semicontinuity and
closure under Borel probability mixtures.  Its proof constructs mixtures
with estimates that hold throughout prescribed neighborhoods, then keeps
all later choices inside those neighborhoods.

\begin{lemma}[Fusion of laws]\label{lem:law-fusion}
Let \(S\) be standard Borel, let \(B\) be compact metric, and let
\(\mathcal M\) be a Borel family of probability measures on \(S\), closed
under Borel probability mixtures.  Suppose a Borel relation
\(\mathcal H\subset S\times B\) defines
\[
 m_z(s)=z\{x\in S:(x,s)\in\mathcal H\}\in[0,1].
\]
Assume that \(\mathcal M\) has a Polish topology generating its Borel sets
for which \((z,s)\mapsto m_z(s)\) is jointly upper semicontinuous.  Finally,
suppose that for every countable \(D\subset B\) and every finite atomless
Borel measure \(\lambda\) on \(B\), some \(z\in\mathcal M\) satisfies
\begin{equation}\label{eq:fusion-countable-tests}
 \|m_z|_D\|_{\ell^2(D)}\le C_0,
 \qquad \int_Bm_z\,\dd\lambda=0.
\end{equation}
Then, for every \(\eps>0\), there exists \(\bar z\in\mathcal M\) with
\(\|m_{\bar z}\|_{\ell^2(B)}\le C_0+\eps\).
\end{lemma}

\begin{proof}
Fix a complete compatible metric on \(\mathcal M\).  The defining
relation and Tonelli's theorem give the mixture identity
\begin{equation}\label{eq:fusion-mixture-identity}
 m_{\int z\,\dd P(z)}(s)=\int m_z(s)\,\dd P(z)
 \quad(s\in B)
\end{equation}
for every Borel probability measure \(P\) on \(\mathcal M\).

\smallskip\noindent
\emph{The small-set ideal.}
We first identify the subsets of candidate laws within which every
prescribed countable and atomless test can still be met.
Call \(A\subset\mathcal M\) small if some countable \(D\) and finite
atomless \(\lambda\) admit no member of \(A\) satisfying
Equation~\eqref{eq:fusion-countable-tests}.  Small sets form a sigma-ideal.
Indeed, for a countable union take the union of the witnessing countable
sets and a strictly positively weighted sum of the witnessing measures,
with weights chosen to give finite total mass.  A law satisfying the two
tests for these combined data would satisfy the tests for each summand.
By hypothesis \(\mathcal M\) is not small.

For a nonsmall closed set \(A\), remove every basis-open set \(V\) for
which \(A\cap V\) is small.  The remaining set \(A'\) is closed and
nonsmall.  Every neighborhood of each point of \(A'\) meets \(A\)
nonsmally.  Only countably many basis-open sets are removed, which is the
reason this conclusion follows from the sigma-ideal property.

\smallskip\noindent
\emph{A finite split.}
Fix a closed nonsmall set $A$, let $A'$ be the set just obtained,
and fix \(r>0\).
An \(r\)-separated tuple in \(B\) has bounded cardinality,
say at most \(N_r\).  Consider the compact family \(T_r\) of measures
\[
 \sum_{j=1}^n w_j\delta_{s_j},
 \qquad 0\le n\le N_r,
 \qquad \dist(s_i,s_j)\ge r\ (i\ne j),
 \qquad w_j\ge0,
 \qquad \sum_jw_j^2\le1.
\]
The empty tuple gives the zero measure.  Its probability barycenters form a
compact convex set \(K_r\) of positive measures on \(B\).  For every
\(\zeta\in K_r\), the weights of its atomic part have \(\ell^2\) norm at
most \(1\): on any fixed finite point set this follows from convexity of
the Euclidean norm, and one then takes the supremum over finite sets.
Write \(\zeta=\zeta_{\rm at}+\lambda\), where \(\lambda\) is atomless.
Nonsmallness of \(A'\), tested against the countable atoms of \(\zeta\)
and this \(\lambda\), gives a point \(z\in A'\) with
\[
 \int m_z\,\dd\zeta
 =\int m_z\,\dd\zeta_{\rm at}\le C_0.
\]

The following finite-mixture step uses the classical minimax strategy
of the Kneser--Fan theorem as presented by Sion
\cite[Theorem~4.2, p.~175]{Sion1958}.  We give the finite-cover and
separation proof here; the surrounding fusion construction is developed
locally in this proof.

For fixed \(z\), the function
\(\zeta\mapsto\int m_z\,\dd\zeta\) is affine and upper
semicontinuous on \(K_r\).  Therefore the open sets
\[
 \left\{\zeta\in K_r:\int m_z\,\dd\zeta<C_0+\eps/4\right\},
 \qquad z\in A',
\]
cover \(K_r\).  Choose a finite subcover with centers
\(z_1,\ldots,z_k\).  Their integral-vector image of \(K_r\) is convex
and misses the open orthant in which every coordinate exceeds
\(C_0+\eps/4\).  Separation from this open convex set gives coefficients
\(\alpha_j\ge0\), \(\sum_j\alpha_j=1\), such that
\[
 \sup_{\zeta\in K_r}\sum_j\alpha_j\int m_{z_j}\,\dd\zeta
 \le C_0+\eps/4.
\]
No closedness of the integral-vector image is needed here.  The separator
has nonnegative coefficients because the orthant is unbounded in every
positive coordinate direction.  Testing on \(T_r\) and using nonnegative
Euclidean dual weights yields
\begin{equation}\label{eq:fusion-finite-split}
 \sup_{F\ r\text{-separated}}
 \left\|\sum_j\alpha_jm_{z_j}|_F\right\|_{\ell^2(F)}
 \le C_0+\eps/4.
\end{equation}

To extend this estimate to neighborhoods of the centers, keep the
coefficients $\alpha_j$ fixed and write its left side as
\[
 \begin{split}
 \Phi_r(y_1,\ldots,y_k)
 &=\sup_{F\ r\text{-separated}}
       \left\|\sum_{j=1}^k\alpha_jm_{y_j}|_F\right\|_{\ell^2(F)}\\
 &=\max_{\substack{0\le n\le N_r,\ s_1,\ldots,s_n\in B\\
                   \dist(s_i,s_l)\ge r\ (i\ne l)\\
                   w_i\ge0,\ \sum_{i=1}^n w_i^2\le1}}
       \sum_{i=1}^n\sum_{j=1}^k w_i\alpha_jm_{y_j}(s_i).
 \end{split}
\]
The case $n=0$ contributes zero.  The configuration and weight spaces in
this maximum form a finite union of compact spaces.  Joint upper
semicontinuity of the hit functions therefore makes $\Phi_r$ upper
semicontinuous jointly in its centers.  The margin in
Equation~\eqref{eq:fusion-finite-split} now gives each center a small
closed neighborhood \(V_j\), of any
prescribed sufficiently small diameter, such that
\begin{equation}\label{eq:fusion-neighborhood-split}
 \left\|\sum_j\alpha_jm_{y_j}|_F\right\|_{\ell^2(F)}<C_0+\eps
\end{equation}
for every simultaneous choice \(y_j\in V_j\) and every \(r\)-separated
finite set \(F\).  Take the neighborhoods to be closed balls contained in
the requisite open neighborhoods.  Each child \(A\cap V_j\) is closed
and nonsmall: the ball $V_j$ contains an open neighborhood of its center,
and every such neighborhood meets $A$ nonsmally.  Discard children with
zero coefficient.

\smallskip\noindent
\emph{Nested splits.}
Starting with \(A=\mathcal M\), perform the split at successive depths
with \(r_n\downarrow0\).  At depth \(n\), require every child to have
diameter at most \(2^{-n}\).  The result is a finitely branching tree of
nonempty nested closed sets, with transition probabilities given by the
coefficients \(\alpha_j\).  Completeness gives a unique point of
\(\mathcal M\) along every infinite branch.  The branch-limit map is
continuous: two branches with the same depth-\(n\) ancestor have limits
in a set of diameter at most \(2^{-n}\).  The transition probabilities
give a Borel probability measure on branches and thus a Borel probability
measure \(P\) on their limits.  By mixture closure,
\(\bar z=\int z\,\dd P(z)\) belongs to \(\mathcal M\).

It remains to check every finite-set estimate for this one law $\bar z$.
Fix a depth \(n\) and an \(r_n\)-separated finite set \(F\).  At a parent
of that depth, let $P_j$ be the distribution of the branch limit
conditioned on taking its $j$th child.  Its transition probability is
$\alpha_j$, and $P_j$ is supported in that child's closed set.  These
conditionings refer to tree nodes, so the closed child sets are allowed
to overlap.  Equation~\eqref{eq:fusion-neighborhood-split} holds for
every simultaneous choice of points from the child sets.  Integrating
with respect to the product of their conditional distributions gives
\[
 \begin{split}
 \left\|\sum_j\alpha_j\int m_y|_F\,\dd P_j(y)\right\|_{\ell^2(F)}
 &\le
 \int\left\|\sum_j\alpha_jm_{y_j}|_F\right\|_{\ell^2(F)}
       \,\dd(P_1\otimes\cdots\otimes P_k)(y_1,\ldots,y_k)\\
 &\le C_0+\eps.
 \end{split}
\]
The first inequality is convexity of the Euclidean norm.  Its left side
is the norm of the conditional mean hit vector at the parent.
Average over the finitely many
parents at this depth, and use
Equation~\eqref{eq:fusion-mixture-identity}, to obtain
\(\|m_{\bar z}|_F\|_2\le C_0+\eps\).
Every finite set of distinct points is \(r_n\)-separated for all
sufficiently large \(n\).  Taking the supremum over finite sets proves
the lemma.
\end{proof}

\subsection{Application of fusion to path laws}

We apply fusion to the family of laws retaining the occupation,
visit-count, and tuple-density bounds from
Lemma~\ref{lem:countable-law}.

\begin{proof}[Proof of Theorem~\ref{thm:positive-law}]
Fix the reference measure \(\tau\) on types and the constant \(C\) from
Lemma~\ref{lem:countable-law}.  Let \(\mathcal M\) be the laws \(z\) on
the base code space \(\mathscr A\) satisfying
\begin{equation}\label{eq:candidate-law-bounds}
 \bar\ell_z=g_z\,\dd a,\quad \|g_z\|_2\le C;
 \qquad \|n_z|_{D_*}\|_{\ell^2(D_*)}\le C+1;
 \qquad \E_z\psi_\Gamma=q_z\tau,\quad \|q_z\|_2\le C.
\end{equation}
We check the structural hypotheses of Lemma~\ref{lem:law-fusion}.

\smallskip\noindent
\emph{Borel structure and mixture closure.}
For a Borel kernel of finite measures on a Polish space with finite reference
measure \(\sigma_0\), the condition that its mean have nonnegative
\(L^2(\sigma_0)\) density of norm at most \(C\) is Borel.  To verify this
directly, choose a countable algebra generating the Borel sets, and test
\[
 \int\!\left(\int h\,\dd\nu_\Gamma\right)\dd z(\Gamma)
 \le C\|h\|_{L^2(\sigma_0)}
\]
for all nonnegative rational simple functions on that algebra, including
the constant function.  These are countably many Borel inequalities.  The
constant test makes the mean measure finite.  The inequalities define a
bounded functional on the dense space of algebra-simple functions in
\(L^2(\sigma_0)\); applying the estimate to absolute values treats signed
functions.  Its representing density gives the mean measure, first on the
algebra and then on all Borel sets.  Thus the tests are also sufficient.

Apply this observation to \(\ell_\Gamma\) and \(\psi_\Gamma\).
The count functions \(N_\Gamma(s)\) are Borel by
Lemma~\ref{lem:path-coding}; finite partial sums of the squared mean counts
test the middle condition of Equation~\eqref{eq:candidate-law-bounds}.
It follows that \(\mathcal M\) is a Borel subset of the standard Borel
space of probabilities on \(\mathscr A\).  The same integral inequalities
and convexity of finite-dimensional Euclidean norms show that all three
conditions are preserved by every Borel probability mixture of members of
\(\mathcal M\).  A mixture is still a law on the base code space, so it
also preserves the endpoint conditions and path-system constraints.

\smallskip\noindent
\emph{Upper semicontinuity of hits.}
Forgetting tags gives a Borel map
\[
 \mathcal M\longrightarrow
 \mathcal P\bigl(C([0,1],X)\bigr).
\]
Choose a Polish topology on \(\mathcal M\), with its same Borel sets,
making this map continuous.  The Borel probability mixtures just checked
remain the same under this refinement.  The relation
\[
 \mathcal H_0=
 \{(\gamma,s)\in C([0,1],X)\times B:\gamma(t)=s
                         \text{ for some }t\in[0,1]\}
\]
is closed: a sequence of witnessing times has a convergent subsequence.
If \(z_j\to z\) in the chosen topology and \(s_j\to s\), the products
of their untagged path laws with \(\delta_{s_j}\) converge weakly to the
corresponding product for \(z,s\).  The closed-set inequality therefore
gives
\[
 \limsup_jm_{z_j}(s_j)\le m_z(s).
\]
The hit functions have the joint upper semicontinuity required by
Lemma~\ref{lem:law-fusion}.  Their defining relation on tagged codes is
Borel by the same closed hit test composed with the forgetful map.

\smallskip\noindent
\emph{The countable tests.}
The structural hypotheses are now established.  The uniform extraction
lemma supplies the remaining test hypothesis within this fixed candidate
family.
Given any countable \(D\subset B\) and finite atomless \(\lambda\), apply
Lemma~\ref{lem:countable-law} to
\(D\cup D_*\cup\{b\}\).  The resulting law belongs to \(\mathcal M\),
and \(m_z\le n_z\) gives
\[
 \|m_z|_D\|_{\ell^2(D)}\le C+1,
 \qquad \int_Bm_z\,\dd\lambda=0.
\]
All hypotheses of Lemma~\ref{lem:law-fusion} hold with \(C_0=C+1\).
Take \(\eps=1\) there.  Its law \(\mu\in\mathcal M\) satisfies
\(\|m_\mu\|_{\ell^2(B)}\le C+2\), while membership in \(\mathcal M\)
retains the occupation and mean-count estimates.  Finally the last
condition in Equation~\eqref{eq:candidate-law-bounds} implies absolute
continuity in each slot and type, by the positive domination in
Equation~\eqref{eq:weighted-tuple-measure}.  This proves every conclusion
with the common bound \(C+2\).
\end{proof}
